% === VIII. DISCUSSION AND LIMITATIONS ========================================
\section{Discussion and Limitations}
\label{sec:disc}

\subsection{What generalises}

The reachable-friction bound~\eqref{eq:mureach} is kinematic. It involves no
tire parameter, no simulator-specific quantity and no property of the residual
model, so it applies to any implementation of the virtual-steady-state-data
strategy of~\cite{ontrack2025}, in simulation or on hardware. Its practical
content is a design rule: \emph{the rollout must be configured to demand at
least the friction one intends to identify}. Where it cannot, $D$ should be
reported as unidentified rather than as fitted. The diagnostic signature,
coefficients resting exactly on box constraints, is likewise generic and should
be checked and logged by any bounded tire-fitting routine, since it is
otherwise indistinguishable in the output from a successful fit.

The two structural properties of Section~\ref{subsec:contraction} generalise on
the same grounds. That stage~3's residual depends on $B$, $C$ and $D$ only
through their product follows from the quasi-steady force recovery
\eqref{eq:ssforces} that defines the virtual-data step, not from any tire or
simulator; any implementation of that step inherits it, and with it the
conclusion that the peak factor must come from the recorded trace rather than
from the rollout. That the loop is not a contraction follows from stage~4
alone, and the remedy, a relaxation on the coefficient update, is one line and
reduces to the inherited behaviour at $\beta=1$. What does not generalise is
the value \num{0.20}, which is a property of this vehicle's extrapolation
error; what generalises is that $\beta=1$ is not defensible and that
$\beta=0$, the undamped loop's opposite, is worse still.

The insufficiency of per-coefficient box constraints for deployment
admissibility is also platform-independent. Both failures we encountered, a
set that satisfies every box constraint while being degenerate and a set that
satisfies every box constraint while being pairwise oversteering and unstable
within the controller's operating range, are properties of the
parameterisation rather than of CARLA. Deep Dynamics~\cite{deepdynamics2024}
constrains coefficients inside the network. We argue that a second,
derived-quantity check belongs at the interface where the model is handed to a
controller.

\subsection{What does not generalise}

Three results are properties of this plant and must not be read as statements
about tires. First, the numerical tire values. The plant is a PhysX
wheeled-vehicle model with a configured friction coefficient, so
Table~\ref{tab:ref} describes \emph{that} plant, and agreement with it is
evidence that the identifier recovers the plant it observes, not evidence
about real tire behaviour. Sagmeister \emph{et al.}~\cite{simfidelity2026} show that
sensitivity to model fidelity is largest near the acceleration limits, which is
where this study operates, and R-CARLA~\cite{rcarla2025} exists because CARLA's
native dynamics are a known weak point for racing-grade work; hardware
replication is therefore necessary before any claim about identification
accuracy on a real vehicle. Second, the gap between configured wheel friction
and realised peak axle friction (\num{1.0} against \num{0.70} on the runs of
Section~\ref{sec:results}, \num{1.5} against \numrange{1.00}{1.05} on the
earlier ones) is an
artefact of the server multiplying the configured wheel coefficient by the road
surface's own, a measured constant \num{0.70} on this map. We report it only
because an identifier, or a reviewer, tuned to converge on the number in the
configuration file will converge on the wrong number, and because the
simulator publishes the effective value per wheel, so nothing forces that
mistake. Third, the \num{0.76} achieved/commanded steering ratio arises from
CARLA's Ackermann controller and this vehicle's steering configuration. The
lesson, identify against the achieved actuator state, is general and standard
practice. The magnitude is not.

\subsection{Limitations}
\label{subsec:limits}

\begin{enumerate}
\item \textbf{One run per configuration, so no metric in
Table~\ref{tab:comparison} carries an uncertainty.} Every value there is a
point estimate from a single \SI{580.4}{\second} and a single
\SI{554.6}{\second} run. The \SI{0.151}{\meter} against \SI{0.224}{\meter} RMS
tracking difference, in particular, has no confidence interval attached and no
evidence that it exceeds run-to-run variation, which is not measured here. The
order-of-magnitude force and friction margins are large enough that run-to-run
variation is unlikely to reverse them, but that is an argument and not a
measurement. Nothing about the experiment prevents repetition, since the sweep
is one command and the scenario file already carries the per-arm parameter
sets, so this is a gap in the evidence rather than in the method.
\item \textbf{Single surface, vehicle and trajectory.} All results come from
one vehicle configuration on one map with the friction held static at
\num{0.70}. That the surface differs from the one the prior was measured on
(Section~\ref{subsubsec:priorprovenance}) makes the identification result a
real one rather than a restatement of the prior, but it is \emph{not} a
demonstration of adaptivity: the identifier is never asked to follow a
\emph{change} in friction. The runtime friction override and the step and decay
schedules the scenario file already carries support that experiment, arguably
the strongest argument for on-track identification, and we do not perform it
here.
\item \textbf{The friction estimator's excitation sweep is synthetic}
(Table~\ref{tab:mu}). The synthetic plant was chosen because $\mu$ is known
there and excitation can be swept independently of it. On the CARLA runs the
warm start is exercised at the one excitation level the lap happens to supply,
where it returns $\mu$ \SI{11}{\percent} high against the \SI{0.2}{\percent} of
the synthetic plant, so the excitation dependence of that bias is not mapped
against the simulator's per-wheel forces. The peak factor the pipeline
publishes is exactly this estimate, so that bias passes through undamped.
\item \textbf{The rear axle is not identified on this trajectory} and the
front-axle $\mu$ is applied to both (Section~\ref{subsec:frictionfloor}). The
quasi-steady argument that makes $\mu_{f}=\mu_{r}$ holds on a homogeneous
surface and would not survive a split-$\mu$ one, which this vehicle's sensor
set cannot detect in the first place.
\item \textbf{The two configurations differ in five settings at once}
(Table~\ref{tab:scenarios}), so Section~\ref{subsec:summary} attributes nothing
to the residual architecture, the loss, the solver, the warm start or the
relaxation individually. Section~\ref{subsec:ablations} does attribute them,
but by offline replay of the identification stage on recorded buffers rather
than in the loop, so its ordering is not a closed-loop ordering. That is also
the status of the S4D accuracy figure there: \SI{21}{\percent} on force RMSE
at matched settings, offline, which we report as a measured ablation and not as
a reproduction of the closed-loop S4 improvement of~\cite{visionsysid2026}. The
excitation-aware rollout, this paper's central contribution, is active in
\emph{both} arms, so the closed-loop comparison of
Section~\ref{subsec:resultsclosedloop} is not evidence for the identifiability
correction; the evidence for that is the rollout-speed sweep of
Section~\ref{subsec:sweep}.
\item \textbf{The physics-informed objective buys nothing measurable here.}
Once its units are reconciled with the data term and the loop is damped, it is
indistinguishable from a plain one-step MSE on this manoeuvre
(Section~\ref{subsec:ablations}). It is retained because it is the paper's
stated formulation and costs nothing, not because it is shown to help.
\item \textbf{Longitudinal dynamics are out of scope.} As in the baseline,
only lateral dynamics are identified. Combined slip and load transfer are
neglected, and the simulator's longitudinal wheel force channel is drivetrain
torque rather than a contact-patch force, so it is unusable as ground truth
for a longitudinal extension without further work.
\end{enumerate}
